% index.tex -- index of Module 8 multiple-choice questions in the 2019-2025 NSW HSC Physics papers.
% Descriptions are paraphrases, not question text.
\documentclass[12pt]{article}
\usepackage{quiz}
\usepackage{longtable}
\hypersetup{pdftitle={HSC Physics Module 8 Multiple-Choice Questions 2019-2025 - Index}}
\begin{document}
\quiztitle{Year 12 Physics}{HSC Module 8 Multiple-Choice Questions 2019–2025: Index}
Every Section I question from the 2019--2025 NSW HSC Physics papers that tests Module 8
(From the Universe to the Atom), with the answer from the NESA marking guidelines. Topics: T1 origins of the elements; T2 structure of the atom; T3 quantum mechanical nature of the atom; T4 properties of the nucleus; T5 deep inside the atom. A second module after
the topic (for example T3/M8) means the question also uses ideas from that module. The descriptions are
short summaries, not the question text.

{\small\renewcommand\arraystretch{1.3}
\begin{longtable}{|c|c|c|p{108mm}|c|}
  \hline
  \textbf{Year} & \textbf{Qu} & \textbf{Topic} & \textbf{What it tests} & \textbf{Ans}\\ \hline
  \endhead
    2019 & 3 & T2 & Geiger-Marsden: particles used and result expected from the plum-pudding model & C\\ \hline
    2019 & 4 & T1 & Compare four stars on an HR diagram (luminosity, colour, temperature, stage) & C\\ \hline
    2019 & 8 & T1 & Velocity-distance graph matching Hubble's original data & A\\ \hline
    2019 & 12 & T5 & Quark change udd to uud: the accompanying product & A\\ \hline
    2019 & 13 & T3/\allowbreak M7 & Photons per second from a laser & B\\ \hline
    2019 & 19 & T4 & Use binding energies to decide whether a reaction releases energy & D\\ \hline
    2020 & 4 & T4 & Decay constant from a mass-time decay graph & B\\ \hline
    2020 & 6 & T1 & Property shared by the stars on a globular-cluster HR diagram & A\\ \hline
    2020 & 8 & T4 & Decay chain on a neutron-proton chart: decay mode and product & C\\ \hline
    2020 & 9 & T3 & Observation supporting Bohr's improvement on Rutherford's model & A\\ \hline
    2020 & 11 & T4 & Classify a neutron-induced reaction from reactant and product masses & C\\ \hline
    2020 & 16 & T4 & Fission reactor losing its moderator: compensating action & A\\ \hline
    2021 & 3 & T5 & Which listed particle is not fundamental in the Standard Model & D\\ \hline
    2021 & 5 & T1/\allowbreak M7 & Identify the spectrum's source and line type & A\\ \hline
    2021 & 6 & T1 & Isotope formed in a step of the proton-proton chain & D\\ \hline
    2021 & 13 & T3 & Bohr-level transitions: which gives the shortest wavelength & B\\ \hline
    2021 & 19 & T4 & Decaying beta source: whether current and stopping voltage change over time & A\\ \hline
    2022 & 3 & T4 & Identify a radiation type from its deflection in an electric field & D\\ \hline
    2022 & 5 & T5 & Quark combination that forms a neutron & B\\ \hline
    2022 & 16 & T4 & Two nuclei's binding energies: separation energy and mass & A\\ \hline
    2023 & 4 & T4 & Mass remaining after three half-lives & B\\ \hline
    2023 & 7 & T3 & De Broglie wavelengths of a proton and a neutron at equal speed & C\\ \hline
    2023 & 11 & T4 & Decay-series chart: two daughters and their decay modes & A\\ \hline
    2023 & 13 & T4 & What a larger total binding energy implies (mass defect) & C\\ \hline
    2023 & 20 & T1 & Velocity-distance graph consistent with accelerating expansion & D\\ \hline
    2024 & 3 & T5 & Which listed particle is fundamental in the Standard Model & C\\ \hline
    2024 & 5 & T1 & Order three star clusters by age from their HR diagrams & A\\ \hline
    2024 & 7 & T4 & Parent : daughter ratio after two half-lives & B\\ \hline
    2024 & 14 & T3 & De Broglie wavelength of an alpha particle against a faster proton & B\\ \hline
    2024 & 19 & T2/\allowbreak M6 & Proton in balanced E and B fields entering at double speed: its acceleration & C\\ \hline
    2025 & 3 & T3 & Whose atomic model explains the hydrogen emission lines & B\\ \hline
    2025 & 4 & T3 & Shortest de Broglie wavelength among electrons and protons at given speeds & D\\ \hline
    2025 & 10 & T4 & Activity-time graphs: valid exponential decay with the smallest decay constant & C\\ \hline
    2025 & 15 & T1 & Two stars of the same class at different magnitudes on an HR diagram & C\\ \hline
    2025 & 16 & T4 & Nucleus that absorbs a neutron then alpha-decays to lithium-7 & A\\ \hline
    2025 & 20 & T5 & Can an electron-positron collision make a proton-antiproton pair, and why & C\\ \hline
\end{longtable}}

Summary: 36 questions, 3 to 6 in each paper (T4: 14, T1: 8, T3: 7, T5: 5, T2: 2).
Properties of the nucleus (half-life, decay charts, binding energy and nuclear reactions) are tested in every paper and most often; HR diagrams, Hubble's law, de Broglie wavelengths, Bohr's model and the Standard Model (quarks and fundamental particles) also come up often.
\end{document}
